\section{Additional results for \texorpdfstring{Section~\ref{sec:ren}}{Section 13}: renormalization}\label{sup:ren}

This section collects statements, examples and remarks of \cref{sec:ren} that the main text cites without stating in full. Proofs are in \cref{app:ren}.

\subsection{Decimation beyond the nearest-neighbor chain}\label{sup:ren-decimation-beyond-the-nearest-neighbor-c}

\begin{proposition}[Nearest-neighbor passing on $\Z$]\label{prop:ren-ring}
Let $W_{ii}=c>0$ and $W_{ij}=a$ for neighbors on $\Z$, with $Q^\zeta=qI$ and $\varepsilon=1-c-2a>0$, so that $\Pi=(I-W^2)/q$ and, by \cref{thm:sca-oz}(e), $\Sigma(x)=\mathsf c_+e^{-m\lvert x\rvert}+\mathsf c_-(-1)^xe^{-m'\lvert x\rvert}$ with $\mathsf c_+=q/(4a\sinh m)$, $\mathsf c_-=q/(4a\sinh m')$, $\cosh m=1+\varepsilon/(2a)$ and $\cosh m'=(1+c)/(2a)$.
\begin{enumerate}
\item[(a)] The two-pole class $\{\mathsf c_+e^{-\mu\lvert x\rvert}+\mathsf c_-\mathsf p^xe^{-\mu'\lvert x\rvert}\}$, $\mathsf p=\pm1$, is closed under decimation by $b$: $(\mu,\mu',\mathsf p)\mapsto(b\mu,b\mu',\mathsf p^b)$, with $\mathsf c_\pm$ fixed.
\item[(b)] After $l\ge1$ steps, with $\mu=b^lm$, $\mu'=b^lm'$ and $\mathsf p=(-1)^{b^l}$, the decimated precision is diffusive and has infinite range: its couplings decay geometrically with ratio $e^{-\arccosh u_\ast}$, where
\[
u_\ast=\frac{\mathsf c_+\sinh\mu\cosh\mu'+\mathsf c_-\sinh\mu'\cosh\mu}{\mathsf p\mathsf c_+\sinh\mu+\mathsf c_-\sinh\mu'}>1.
\]
\item[(c)] For $\lvert x\rvert\ge1$ the staggered term after $l$ steps is at most $(\sinh m/\sinh m')e^{-b^l(m'-m)}$ times the smooth one, so the coarse covariance converges in relative terms to that of a nearest-neighbor precision with mass $b^lm$, plus a nugget $\mathsf c_-\one_{\{x=0\}}$ of relative size $\sinh m/\sinh m'<1$, which vanishes as $\varepsilon\to0$.
\end{enumerate}
\end{proposition}
% src: H4-P2 (CHECKED), with the analytic step for odd b inserted per HC_critic V8(i)
Proof in \cref{app:ren}, including the step for odd $b$, where the staggered pole changes sign; checked in \cref{chk:ren-ring}. The range-two structure and the staggered pole are irrelevant: the medium flows into the class of \cref{thm:ren-decimation}. In two dimensions finite-range classes do not close.

\begin{example}[Nearest-neighbor precisions do not close on $\Z^2$]\label{ex:ren-two-dim}
Let $\Pi$ be the nearest-neighbor precision on $\Z^2$, $\Pi_{ii}=\varepsilon^\Pi+4a^\Pi$. Decimating the sites with $x_1+x_2$ odd leaves the even sublattice, a copy of $\Z^2$ in the coordinates $((x_1+x_2)/2,(x_1-x_2)/2)$, with nearest-neighbor couplings $2(a^\Pi)^2/(\varepsilon^\Pi+4a^\Pi)$, diagonal couplings $(a^\Pi)^2/(\varepsilon^\Pi+4a^\Pi)$ and leak $\varepsilon^\Pi(\varepsilon^\Pi+8a^\Pi)/(\varepsilon^\Pi+4a^\Pi)$, so the nearest-neighbor class is left after one step. On a $16\times16$ torus with $\varepsilon^\Pi=0.05$ and $a^\Pi=1$, a second decimation, onto $2\Z\times2\Z$, gives nonzero positive couplings at all $63$ displacements of the coarse torus.
\end{example}
% src: H4-X1 (counterexample, CHECKED)
The first part counts common neighbors (\cref{app:ren}), the second is a computation (\cref{chk:ren-two-dim}).

\begin{observation}[Decay of decimated kernels]\label{obs:ren-decay}
For the nearest-neighbor precision on $\Z^2$ with $a^\Pi=1$, decimated onto $2\Z\times2\Z$ on a $256\times256$ torus at $\varepsilon^\Pi=0.5$, $0.05$, $0.005$ and $0.0005$, the coarse precision has no positive off-diagonal entry, and at coarse displacements up to five its axis couplings fall by factors between $0.025$ and $0.095$ per coarse step, between $0.090$ and $0.095$ from the second step on at the three smaller leaks; onto $4\Z\times4\Z$ the factors at displacements up to five lie between $0.019$ and $0.129$. The factor rises slowly with distance: at $\varepsilon^\Pi=0.0005$ the ratio of the couplings at coarse displacements nine and eight is $0.096$ onto $2\Z\times2\Z$ and $0.135$ onto $4\Z\times4\Z$. The decay per coarse step changes little between $\varepsilon^\Pi=0.05$ and $0.0005$, by at most a quarter at each step; at $\varepsilon^\Pi=0.5$, a leak comparable to the coupling, it is faster, and onto $4\Z\times4\Z$ the factors there, at displacements up to five, lie between $0.019$ and $0.030$. We have not proved this; the eliminated sites form a random walk killed on the retained lattice, which suggests a rate that does not depend on the leak when the leak is small against the coupling.
\end{observation}
% src: H4-X1 observation (open item O8); recomputed at four leaks by ren_checks.py part A

\subsection{Block averages, exponents and time}\label{sup:ren-block-averages-exponents-and-time}

\begin{example}[Block averages leave the diffusive class]\label{ex:ren-block-average}
Let $\Sigma(x)=\mathsf c_\Sigma e^{-m\lvert x\rvert}$, $\mathsf c_\Sigma=1/(2a^\Pi\sinh m)$, be the nearest-neighbor field on $\Z$, and let $b\ge2$. Then
\[
\Sigma^{\mathsf{Avg}}(X)=\frac{\mathsf c_\Sigma\mathsf g_b}{b^2}\bigl(e^{-bm\lvert X\rvert}-\mathsf k_b\delta_{X0}\bigr),\qquad\mathsf g_b=\frac{\sinh^2(bm/2)}{\sinh^2(m/2)},
\]
with $\mathsf k_b=1-\mathsf g_b^{-1}(b+2\sum_{j=1}^{b-1}(b-j)e^{-jm})\in(0,1)$, the covariance of an autoregressive moving-average field of order $(1,1)$. The precision of the block averages has the entries $\Pi^{\mathsf{Avg}}(n)=-\mathsf c_1(-\mathsf h_b)^{\lvert n\rvert-1}$ for $\lvert n\rvert\ge1$, with $\mathsf c_1>0$, $\mathsf h_b=(\mathsf e_0-(\mathsf e_0^2-\mathsf k_b^2)^{1/2})/\mathsf k_b\in(0,1)$ and $\mathsf e_0=\sinh bm-\mathsf k_b\cosh bm$, so it has positive entries, that is negative couplings, at every even distance, for every $\varepsilon^\Pi>0$ and $b\ge2$. As $m\to0$, $\mathsf h_b\to(1-(1-\mathsf t_b^2)^{1/2})/\mathsf t_b$ with $\mathsf t_b=(b^2-1)/(2b^2+1)$, which is $0.1716$ for $b=2$ and tends to $2-\sqrt3=0.2679$ as $b\to\infty$: the defect does not vanish near the massless point.
\end{example}
% src: H4-X2 (counterexample, CHECKED); C3_04 recomputation
This form is the spatial counterpart of the temporal aggregation of an autoregression \citep{amemiya1972effect,working1960note}, works cited from their records for attribution only. Proof in \cref{app:ren}; checked in one and two dimensions in \cref{chk:ren-block-average}. The exact law of the natural meta-agent variables thus leaves the class on which contrast is defined (\cref{open:col-non-diffusive}), even for reversible diffusive populations; the error $\mathcal D$ of the coarse medium is small for blocks with $\Vert\Pi_{\perp\Phi}\Vert^2\ll\xin^{\min}\xout$, which high contrast alone does not ensure (\cref{thm:ren-memory}(d), \cref{sec:ren-closure}), and is not small without contrast (\cref{ex:ren-no-contrast}). The size bound of \cref{thm:col-contrast-size} concerns diffusive precisions and does not apply to the covariance-side length of block averages in \cref{thm:ren-flows}(c).

\begin{proposition}[Fixed points, exponents and time]\label{prop:ren-exponents}
Let $\Pi$ be translation invariant and diffusive on $\Z^D$ with uniform leak $\varepsilon^\Pi>0$, and let $\hat\Pi-\varepsilon^\Pi$ be positive away from $2\pi\Z^D$.
\begin{enumerate}
\item[(a)] If the couplings have finite second moments, decimation and block averaging send $\mathfrak x_{2,d}=1/\xi_{2,d}^2$ to $b^2\mathfrak x_{2,d}(1+O(b^2\mathfrak x_{2,d}))$: near the massless point the leak is relevant with $y_{\mathrm{RG}}=2$. On a finite closed population \textup{(M)} forces a positive mean leak, so the massless point is never reached.
\item[(b)] (Long-range couplings.) Let instead $\hat\Pi(p)=\varepsilon^\Pi+\mathsf a_\sigma\lvert p\rvert^\sigma+O(\lvert p\rvert^{\sigma'})$ near $p=0$, with $0<\sigma<\sigma'\le2$ and $\mathsf a_\sigma>0$, and let $\hat\Pi$ be twice continuously differentiable away from $2\pi\Z^D$. Then the coarse symbols have the form $\varepsilon^{\Pi,(1)}+\mathsf a^{(1)}_\sigma\lvert p'\rvert^\sigma+O(\lvert p'\rvert^{\min(\sigma',2\sigma,2)})$, with $\varepsilon^{\Pi,(1)}=b^D\varepsilon^\Pi/(1+\varepsilon^\Pi\mathsf S_1)$ and $\mathsf a^{(1)}_\sigma=b^{D-\sigma}\mathsf a_\sigma/(1+\varepsilon^\Pi\mathsf S_1)^2$ for decimation and $\varepsilon^{\Pi,(1)}=b^D\varepsilon^\Pi$ and $\mathsf a^{(1)}_\sigma=b^{D-\sigma}\mathsf a_\sigma$ for block averaging. The ratio $\varepsilon^\Pi/\mathsf a_\sigma$ is relevant with $y_{\mathrm{RG}}=\sigma$, the couplings have infinite second moment, $\xi_2=\infty$, and the size bound of \cref{cor:col-lstar} is void.
\item[(c)] (Time.) Under the reversible diffusion with drift $\Pi$, each term $k$ of the decimated autocovariance relaxes at $\hat\Pi((p'+2\pi k)/b)$. The term $k=0$ relaxes at $\varepsilon^\Pi+\sum_xa^\Pi(x)(p'\cdot x)^2/(2b^2)+o(\lvert p'\rvert^2b^{-2})$, which is $\varepsilon^\Pi+\sum_d\varsigma_d^2p_d'^2/(2b^2)+o(\lvert p'\rvert^2b^{-2})$ when the mixed moments $\sum_xa^\Pi(x)x_dx_{d'}$, $d\neq d'$, vanish, so in the time unit $b^2$ the coupling is invariant and the leak is multiplied by $b^2$ (dynamic exponent $2$, and $\sigma$ under (b)); for the nearest-neighbor kernel the other terms relax at rates at least $\varepsilon^\Pi+4a^\Pi\sin^2(\pi/(2b))$. For the level-0 recursion one coarse round is $b^2$ base rounds.
\end{enumerate}
\end{proposition}
% src: H4-P4 (1), (3); (2) restated with explicit hypotheses and proved, checked by ren_checks.py part C; quartic clause moved to rem:ren-quartic per HC_critic V8(ii); finite-population clause per decision 16
Proof in \cref{app:ren}; part (b) is checked for the kernel $1/x^2$ on $\Z$ in \cref{chk:ren-exponents}. The exponent $y_{\mathrm{RG}}=\sigma$ is the thermal eigenvalue of the Gaussian fixed point for long-range couplings in the renormalization-group treatment of \citet{fisher1972critical}, a work cited from its record and abstract for attribution, where that fixed point governs the interacting model for $\sigma<D/2$; part (b) restates it for decimation and block averaging of a lattice precision with explicit remainder terms, and for the Gaussian fields considered here it holds for every $0<\sigma<2$. With $\tau\simeq\xi_2^2/a$ (\cref{prop:sca-diffusive-scaling}(a)), the coarse length in block units and the coarse correlation time in units of $b^{2l}$ base rounds reach order one at the same level: space and time terminate together.

\begin{remark}[The quartic moment]\label{rem:ren-quartic}
At the massless point the term $k=0$ of \eqref{eq:ren-symbols} multiplies the ratio of the quartic to the quadratic coefficient of $\hat\Pi$ by $b^{-2}$, which alone would make the shape of the kernel beyond its second moment irrelevant with $\Lambda_{\mathrm{RG}}=b^{-2}$. The aliased terms contribute to the quartic coefficient at the same order: in one dimension, with $\hat\Pi(p)=\alpha_2p^2+\alpha_4p^4+O(p^6)$ at $\varepsilon^\Pi=0$, decimation gives the ratio $b^{-2}(\alpha_4/\alpha_2-\alpha_2\mathsf S^{(0)}_1)$, so the term $k=0$ does not decide the flow of the shape. For the nearest-neighbor kernel the two contributions keep the ratio at its lattice value $-1/12$, as the closure of \cref{thm:ren-decimation} requires.
\end{remark}
% src: H4-P4(1) quartic clause, as a leading-order remark per HC_critic V8(ii); aliased contribution derived while drafting

\subsection{Collectives across levels}\label{sup:ren-collectives-across-levels}

\begin{proposition}[The exchangeable comparison admits only dense collectives]\label{prop:ren-exchangeable}
Let a finite closed population of $N$ agents with no declared geometry or coupling mechanism have precision $\varepsilon^\Pi I+L^\Pi$ with uniform precision leak $\varepsilon^\Pi>0$ and mean total coupling $\bar d$ per agent, so that its exchangeable comparison medium has coupling $\bar a=\bar d/(N-1)$ between every pair. Then $\bar r_{\mathrm{hom}}=1+N\bar a/\varepsilon^\Pi$, and a region $C$ of $n\ge2$ agents is a collective at some admissible level if and only if
\begin{equation}\label{eq:ren-exchangeable}
\lambda_2(L_C)-\frac{\cut(C)}n>N\bar a\qquad\text{and}\qquad\varepsilon^\Pi>\frac{N\bar a\cut(C)/n}{\lambda_2(L_C)-\cut(C)/n-N\bar a}.
\end{equation}
Since $\lambda_2(L_C)\le\frac n{n-1}\min_{i\in C}d_C(i)$, with $d_C(i)$ the coupling of member $i$ to the other members, every member of a collective has $d_C(i)>\frac{(n-1)N}{n(N-1)}\bar d$: collectives are denser than the population on average.
\end{proposition}
% src: H4-P7 restated as the consequence of the exchangeable clause (H6-P11); consequence (ii) replaced per H6-P11
Proof in \cref{app:ren}. Had the bridge population of \cref{sec:col-bridge} no declared geometry, $N\bar a=3.382$ would exceed the gap $0.268$ of each grid, and neither grid nor their union would ever be a collective; with its declared $6\times12$ box each grid is a collective on an interval of leaks, and the union, the whole closed population, never is (\cref{prop:col-whole-population}). Structure contained in a declared geometry is assumed rather than found (\cref{rem:col-declared-geometry}).

\begin{proposition}[Nested collectives judged level by level]\label{prop:ren-per-level}
In the nested medium of \cref{thm:ren-nested}, with the exchangeable comparison at every level and with $\Delta_{\mathrm H}$ and $T_{\mathrm H}$ as defined in \cref{sec:ren-levels}:
\begin{enumerate}
\item[(a)] judged in the level-0 medium, blocks of level two or more are never collectives, and level-1 blocks are collectives at some admissible level if and only if $\Delta_{\mathrm H}(b,\mathfrak q,n_{\mathrm H})>0$ and $\varepsilon^\Pi>A_{\mathrm h}T_{\mathrm H}(b,\mathfrak q,n_{\mathrm H})$;
\item[(b)] for $0\le l\le n_{\mathrm H}-2$, the level-$(l+1)$ blocks are collectives of the level-$l$ coarse medium if and only if $\Delta_{\mathrm H}(b,\mathfrak q,n_{\mathrm H}-l)>0$ and $\varepsilon^\Pi/(A_{\mathrm h}\mathfrak q^l)>T_{\mathrm H}(b,\mathfrak q,n_{\mathrm H}-l)$, and the whole population is never a collective of the level-$(n_{\mathrm H}-1)$ medium;
\item[(c)] for $n_{\mathrm H}\ge2$, $\Delta_{\mathrm H}(b,\mathfrak q,m)>0$ for every $m=2,\dots,n_{\mathrm H}$ if and only if $\mathfrak q<b/(2b^2-1)$, so nested collectives exist at every level $1,\dots,n_{\mathrm H}-1$, judged level by level, if and only if $\mathfrak q<b/(2b^2-1)$ and $\varepsilon^\Pi>A_{\mathrm h}\max_{0\le l\le n_{\mathrm H}-2}\mathfrak q^lT_{\mathrm H}(b,\mathfrak q,n_{\mathrm H}-l)$;
\item[(d)] if $\mathfrak q\ge1/b$, no block of any level is a collective at any leak, under either judgment; this contains the regime $1/b<\mathfrak q<1$ in which \cref{prop:abs-hierarchy}(b) gives power-law coherence;
\item[(e)] if $\mathfrak q<1/b$ and $k<l$, the level-$(l+1)$ blocks have contrast below the baseline of the level-$k$ medium, so a single level declared for all levels admits collectives at one level at most;
\item[(f)] if $\mathfrak q\ge1/(2b-1)$, then $\Delta_{\mathrm H}(b,\mathfrak q,m)<0$ for every $m\ge2$, and the conclusion of (d) holds; if $\mathfrak q<1/(2b-1)$, then $\Delta_{\mathrm H}(b,\mathfrak q,m)$ increases strictly with $m\ge2$ to $(2b-1)/(b-1)-2/(1-\mathfrak q)>0$, so with $m^\ast$ the least $m\ge2$ at which it is positive, which is $2$ exactly when $\mathfrak q<b/(2b^2-1)$, the level-$(l+1)$ blocks are collectives of the level-$l$ coarse medium at some leak if and only if $l\le n_{\mathrm H}-m^\ast$, and then at every leak above the threshold of (b).
\end{enumerate}
\end{proposition}
% src: H6-T2 (CHECKED, H6_nested.py, H6_misc.py part 2); replaces H4-R1 and H4-T18(6)(iv) per H6-P13; (f) from the figure report on fig_ren_per_level, proved and checked in the verification round (pvrf_ren)
Proof in \cref{app:ren}; checked in \cref{chk:ren-per-level}, where the medium $(b,n_{\mathrm H},A_{\mathrm h},\mathfrak q)=(2,8,1,0.2)$ has collectives at every level for $\varepsilon^\Pi>0.634$. Block contrasts are the same in every coarse medium (\cref{thm:ren-nested}(c)), while the excess of the exchangeable baseline over one falls like $\mathfrak q^l$; hence the dependence on $\varepsilon^\Pi/(A_{\mathrm h}\mathfrak q^l)$ and the need for a declared level at each level. Without declared geometry, scale-free nested coherence and nested collectives exclude each other, and part (f) extends the exclusion below $1/b$: for $b=2$ the blocks of every level qualify at large enough leaks when $\mathfrak q<2/7$; when $2/7\le\mathfrak q<1/3$ the blocks of the top $m^\ast-2$ levels below the whole population never qualify while all lower levels do at large enough leaks, with $m^\ast=3$ at $\mathfrak q=0.3$ and $m^\ast=7$ at $\mathfrak q=0.33$; and when $\mathfrak q\ge1/3$ no level qualifies. A declared geometry admits more: with the sixteen agents of the medium with $b=2$, $n_{\mathrm H}=4$, $A_{\mathrm h}=1$ at the sites $0,\dots,15$ of $\Z$ and the lattice clause at every level, exhaustive enumeration finds collectives at all three levels for $\mathfrak q\le0.55$ at strong enough grounding; the general embedding is open. \Cref{fig:ren-per-level} places the conditions of \cref{prop:ren-per-level}(c), (d) and (f) in the plane of $b$ and $\mathfrak q$ and shows, for the medium $(b,n_{\mathrm H},A_{\mathrm h},\mathfrak q)=(2,8,1,0.2)$, how the leak thresholds fall like $\mathfrak q^l$ from $0.6336$ at $l=0$.

\subsection{Memory without contrast and the elimination of internal states}\label{sup:ren-memory-without-contrast-and-the-eliminat}

\begin{example}[Without contrast, memory is the leading effect]\label{ex:ren-no-contrast}
In a homogeneous medium the Galerkin description, equivalently the Markov coarse network obtained by averaging the update weights, overstates the coarse coupling by the block factor. For the nearest-neighbor precision on $\Z$ and blocks of $b$ sites, the coarse medium has coupling $a^\Pi$ between adjacent blocks and leak $b\varepsilon^\Pi$, so its squared second-moment length in block units is $a^\Pi/(b\varepsilon^\Pi)$, while the block averages have $a^\Pi/(b^2\varepsilon^\Pi)+O(1)$ (\cref{thm:ren-flows}(c)): the ratio tends to $b$. For the level-0 recursion with $W_{ii}=c$ and $W_{ij}=a$ between neighbors on $\Z$, averaging the weights over blocks gives the coarse message weight $a/b$ per round, while a nearest-neighbor coarse network with the same leak per base round reproduces the exact axis rate $bm$ of \cref{thm:ren-flows}(b) only with the message weight $\varepsilon/(2(\cosh bm-1))=(a/b^2)(1+O(\varepsilon))$.
\end{example}
% src: H4-X3 (CHECKED)
Proof in \cref{app:ren}; checked in one and two dimensions in \cref{chk:ren-no-contrast}. Within a block without contrast, internal disagreement relaxes no faster than content crosses the block, so treating each block as internally instantaneous transports content $b$ times too fast.

The grouping step $G$ of the adiabatic reduction described in \cref{sec:ren-memory} eliminates internal states. Let $x=(x_{\mathrm{ret}},x_{\mathrm{int}})$ collect the states of a population of particles, $x_{\mathrm{int}}$ their internal states and $x_{\mathrm{ret}}$ all others, evolving as $\mathrm dx=-Mx\mathrm dt+\sqrt2\mathsf N\mathrm d\mathsf B_t$ with $\mathsf N$ block diagonal over the four blocks of the particular partition, as in \citet{friston2023path}, and every eigenvalue of $M$ with positive real part. Write $M_{\mathrm{rr}},M_{\mathrm{ri}},M_{\mathrm{ir}},M_{\mathrm{ii}}$ and $\mathsf N_{\mathrm r},\mathsf N_{\mathrm i}$ for the blocks, $\Sigma$ for the stationary covariance, $\mathsf C_{\mathrm{ir}}=\Sigma_{\mathrm{ir}}\Sigma_{\mathrm{rr}}^{-1}$ for the regression of internal on retained states, and $x_\perp=x_{\mathrm{int}}-\mathsf C_{\mathrm{ir}}x_{\mathrm{ret}}$. In the linear-Gaussian case $\mathsf C_{\mathrm{ir}}x_{\mathrm{ret}}$ is the conditional expectation that the elimination substitutes (\citet{friston2019free} conditions on the particle's own blanket and asserts that the two coincide), and $x_\perp$ is what that work calls the intrinsic fluctuations.

\begin{proposition}[What the elimination of internal states drops]\label{prop:ren-adiabatic}
In this setting:
\begin{enumerate}
\item[(a)] $\mathrm dx_{\mathrm{ret}}=-(M_{\mathrm{rr}}+M_{\mathrm{ri}}\mathsf C_{\mathrm{ir}})x_{\mathrm{ret}}\mathrm dt-M_{\mathrm{ri}}x_\perp\mathrm dt+\sqrt2\mathsf N_{\mathrm r}\mathrm d\mathsf B_{\mathrm r}$, and at stationarity $\Cov(x_\perp(t),x_{\mathrm{ret}}(t))=0$;
\item[(b)] the residual is a state, not noise: $\mathrm dx_\perp=-M_{\perp\mathrm r}x_{\mathrm{ret}}\mathrm dt-M_{\perp\perp}x_\perp\mathrm dt+\sqrt2(\mathsf N_{\mathrm i}\mathrm d\mathsf B_{\mathrm i}-\mathsf C_{\mathrm{ir}}\mathsf N_{\mathrm r}\mathrm d\mathsf B_{\mathrm r})$, with $M_{\perp\mathrm r}=M_{\mathrm{ir}}+M_{\mathrm{ii}}\mathsf C_{\mathrm{ir}}-\mathsf C_{\mathrm{ir}}M_{\mathrm{rr}}-\mathsf C_{\mathrm{ir}}M_{\mathrm{ri}}\mathsf C_{\mathrm{ir}}$ and $M_{\perp\perp}=M_{\mathrm{ii}}-\mathsf C_{\mathrm{ir}}M_{\mathrm{ri}}$, so the retained states obey exactly
\[
\dot x_{\mathrm{ret}}(t)=-(M_{\mathrm{rr}}+M_{\mathrm{ri}}\mathsf C_{\mathrm{ir}})x_{\mathrm{ret}}(t)+\int_0^t\Gamma_{\mathrm F}(t-t')x_{\mathrm{ret}}(t')\mathrm dt'+\mathfrak f(t),\qquad\Gamma_{\mathrm F}(u)=M_{\mathrm{ri}}e^{-M_{\perp\perp}u}M_{\perp\mathrm r},
\]
the equation of \cref{thm:agy-mori-zwanzig}(b) for the Mori complement $x_\perp$, where $\mathfrak f$ collects the white noise of $x_{\mathrm{ret}}$, the initial-state term $-M_{\mathrm{ri}}e^{-M_{\perp\perp}t}x_\perp(0)$ and a colored force; the elimination keeps the first term and replaces $-M_{\mathrm{ri}}x_\perp$ by white noise independent of the past of $x_{\mathrm{ret}}$, so it drops the kernel $\Gamma_{\mathrm F}$, the initial-state term and the color of the force;
\item[(c)] the retained process is Markov with a transition law that does not depend on the initial law if and only if $M_{\mathrm{ri}}=0$, which fails for every particle whose active states read its internal states, the particles the construction is built for.
\end{enumerate}
\end{proposition}
% src: H1-P1(a)-(c) (CHECKED); C3_05 recomputation; H1's regression matrix renamed C_ir (the gain symbol is reserved)
Proof in \cref{app:ren}; checked in \cref{chk:ren-adiabatic}. The size of what is dropped is observed, not proved.

\begin{observation}[The error of the memoryless reduction]\label{obs:ren-adiabatic-rate}
In a random instance with three external, two sensory, two active and three internal states, with the internal rows of $M$ and their noise intensity multiplied by a speed factor $\mathsf s$, the largest deviation over lags $u\in[0,3]$ between the exact stationary autocovariance of the retained states and that of the memoryless reduction, with white noise chosen to reproduce $\Sigma_{\mathrm{rr}}$, relative to $\Vert\Sigma_{\mathrm{rr}}\Vert$, is $2.6\times10^{-2}$, $2.2\times10^{-2}$, $1.0\times10^{-2}$, $3.7\times10^{-3}$ and $1.2\times10^{-3}$ at $\mathsf s=1$, $3$, $10$, $30$ and $100$, and a second instance drawn independently behaves alike (\cref{chk:ren-adiabatic}). From $\mathsf s=10$ on the error falls like $1/\mathsf s$, first order in the ratio of timescales, as the expansions in that ratio \citep{vankampen1985elimination,geigenmuller1983systematic}, known to us from their abstracts, lead one to expect; we have not proved a rate.
\end{observation}
% src: H1-P1(d) (numerical observation per HC_critic V10); H1_checks check A; C3_05 independent instance; speed factor renamed s
Three further steps add errors that \cref{prop:ren-adiabatic} does not cover: the asserted independence of one particle's internal states from the blankets of others is a stationary blanket condition, which fails at second order in the coupling for non-normal drift (\cref{sec:rel-debate}); the slow eigenmodes of each blanket are a second projection, with its own memory (\cref{prop:ren-perron}); and correlations among the induced fluctuations are set aside. In a path formulation the exact replacement for the expected internal state given the blanket is the filter of \cref{thm:agy-closure-identity}(d), of which the elimination is the zero-memory truncation.

\begin{remark}[Renormalizing generative models]\label{rem:ren-rgm}
The renormalizing generative models of \citet{friston2025pixels}, read in their primary text, are discrete hierarchical models in which a state at one level generates the initial states and paths of a group of states below; groups do not overlap, so the latent factors of a level are conditionally independent given the level below by construction. Blocking is by tessellation, with a singular value decomposition of each local group (the main text keeps $32$ principal vectors of $4\times4$ pixel groups, the caption of its Figure~4 at most $16$) followed by quantization; higher levels encode sequences of sequences, which, as the authors observe, removes the Markov property of the lowest-level content. Applied to one neuroimaging dataset, $G$ and $R$ gave time constants growing by about $2.37$ and spatial supports by about $2.13$ per scale \citep{friston2021parcels}, figures read through a targeted extraction of the full text. In the parts of the two texts that we read we found no error bound for the reduction. We have not read the later treatment of these models by \citet{zaghw2026renormalising}. In our interpretation, renormalizability in these models is built into the generative architecture rather than derived from a given microscopic dynamics, whereas in a given linear-Gaussian network the retained modes carry the static error $\mathcal D$ and the memory $\Gamma$ (\cref{thm:ren-memory}), both zero exactly when the precision leaves their span invariant; their temporal hierarchy of paths plays the role of the filter state of \cref{thm:agy-closure-identity}(d), a Markov augmentation that absorbs memory.
\end{remark}
% src: H4-R3 (literature, P, with interpretation); H1-R3 (P-qfe, factors unverified by the proof check); HC_critic Section 4 (32 in the main text, 16 in Figure 4's caption)

\subsection{Codes, ranks and interpretation across levels}\label{sec:ren-codes}

A meta-agent $J$ that speaks through its members can aggregate their public messages, $Y_J=\sum_{j\in J}\mathsf w_jy_j$, with weights $\mathsf w_j\ge0$ that sum to one and the posted messages $y_j=H_jm_j+\nu_j$, where $H_j=T_cU_j$, $m_j=\eta+e_j$ and the channel noises $\nu_j\sim\Ncal(0,R_c)$ are independent.

\begin{proposition}[Aggregated codes]\label{prop:ren-codes}
Write $\bar U_J=\sum_j\mathsf w_jU_j$, $H^{\mathrm{agg}}_J=T_c\bar U_J$, $\bar e_J=\sum_j\mathsf w_je_j$ and $\tilde e_j=e_j-\bar e_J$. Then
\[
Y_J=H^{\mathrm{agg}}_J(\eta+\bar e_J)+\sum_{j\in J}\mathsf w_j(H_j-H^{\mathrm{agg}}_J)\tilde e_j+\sum_{j\in J}\mathsf w_j\nu_j.
\]
\begin{enumerate}
\item[(a)] $\rank H^{\mathrm{agg}}_J\le r_c$: aggregation in public coordinates never raises the rank of a code, while concatenating the messages gives a code with row space $\sum_j\operatorname{row}(H_j)$, of rank up to $\min(K,\lvert J\rvert r_c)$.
\item[(b)] With a common effective code the middle term vanishes and the aggregate reads the collective belief $\eta+\bar e_J$ through that code. With diverse codes, uniform weights, uniform leak and a world-conditional covariance of the form $\Sigma^{\mathrm{sc}}\otimes I_K$, the world-conditional variance of the middle term, $\E\lvert\cdot-\E[\cdot\mid\eta]\rvert^2$, is at most $\max_j\Vert H_j-H^{\mathrm{agg}}_J\Vert^2K(n-1)/(n\xin(J))$, with $n=\lvert J\rvert$: contrast suppresses the disagreement term.
\item[(c)] A level-0 meta-receiver of $Y_J$, which assigns the noise $\sum_j\mathsf w_j^2R_c$, ignores $\bar e_J$: for a scalar world, a common code $H$ and uniform weights, the second moment of $Y_J-H\eta$ given the world path is at least $H^2/(n\xout(J))+R_c/n$, so the believed precision exceeds the true one by at least the factor $1+H^2/(\xout(J)R_c)$, the more so the slower the collective's outward rate (\cref{sec:mis}).
\end{enumerate}
\end{proposition}
% src: H4-P11 (DERIVED); level-0 overvaluation per HC_critic 2.8 (decision 10)
Proof in \cref{app:ren}. Across levels, aggregation averages frames.

\begin{proposition}[Aggregated codes across levels]\label{prop:ren-code-levels}
Let the members' frames $U_j$ be independent draws of a law on $\GL(K)$ with finite second moment, and let each level average $b$ frames of the level below, so that the level-$l$ aggregate $\bar U^{(l)}$ averages $b^l$ independent frames. Then $\bar U^{(l)}\to\E U$ almost surely, $\E\Vert\bar U^{(l)}-\E U\Vert_F^2=b^{-l}\E\Vert U-\E U\Vert_F^2$, and the singular values of $T_c\bar U^{(l)}$ converge to those of $T_c\E U$.
\begin{enumerate}
\item[(a)] Frames drawn from the Haar law on $\mathrm{SO}(K)$ or $\mathrm O(K)$, $K\ge2$, have $\E U=0$: the aggregate code vanishes, with $\E\Vert T_c\bar U^{(l)}\Vert_F^2=b^{-l}\Vert T_c\Vert_F^2$.
\item[(b)] Rotations of $\R^3$ about $\mathbf e_1$ by angles $\vartheta$ with a symmetric law give $\E U=\diag(1,\E\cos\vartheta,\E\cos\vartheta)$; a code reading $(\mathbf e_1,\mathbf e_2)$ keeps the singular values $(1,\lvert\E\cos\vartheta\rvert)$, so its effective rank at a relative threshold $\mathsf f_{\mathrm{th}}$ drops from two to one exactly when $\lvert\E\cos\vartheta\rvert<\mathsf f_{\mathrm{th}}$.
\item[(c)] Rotations within the plane that the code reads attenuate both directions equally and keep the rank.
\item[(d)] Receivers that carry messages into world coordinates with true translations (\cref{sec:fw}) before aggregating aggregate identity frames and lose nothing.
\end{enumerate}
\end{proposition}
% src: H4-P12 (CHECKED)
Proof in \cref{app:ren}; checked in \cref{chk:ren-codes}. Aggregate rank falls with scale only in directions on which members' frames disagree, and only under aggregation in public coordinates. Relays of meta-level messages have a rank of their own. In a common prior metric noiseless hops are orthogonal projections (\cref{thm:sca-alternating}); let a meta-level hop apply a convex combination $\bar P=\sum_j\mathsf w_jP_{V_j}$, $\mathsf w_j>0$, as when each member of the receiving meta-agent relays one member of the sending one and the content is averaged.

\begin{proposition}[Meta-level relays]\label{prop:ren-relay-rank}
\begin{enumerate}
\item[(a)] $\bar P$ is symmetric with $0\preceq\bar P\preceq I$, $\bar Px=x$ if and only if $x\in V_\cap=\bigcap_jV_j$, and $\rank\bar P=\dim\sum_jV_j\ge\max_j\dim V_j$.
\item[(b)] Along a stationary relay of meta-level hops, $\sigma_k(\bar P^n)/\sigma_1(\bar P^n)=(\lambda_k(\bar P)/\lambda_1(\bar P))^n$, so the effective rank at any relative threshold is nonincreasing in $n$, and $\Vert\bar P^n-P_\cap\Vert=\lambda_\ast^n$, with $\lambda_\ast<1$ the largest eigenvalue of $\bar P$ on $V_\cap^\perp$.
\end{enumerate}
Hence \cref{conj:sca-rank-down} holds for stationary averaged relays in a common metric, with limiting rank $\dim V_\cap$ when $V_\cap\neq\{0\}$, while in the reading that meta-agents communicate at lower rank than their members it fails at one hop, because averaging parallel paths raises the rank before relay depth lowers it.
\end{proposition}
% src: H4-P13 (CHECKED)
Proof in \cref{app:ren}; in \cref{chk:ren-relay} twelve rank-two codes in $\R^4$ that share one direction give $\bar P$ of rank four and effective ranks $4,3,1,1,\dots$ along the relay. Decisions aggregate as codes do.

\begin{proposition}[The Hayek dimension under aggregation]\label{prop:ren-hayek}
In the setting of \cref{prop:abs-hayek}, let units $C$ of level $l$ have decision functionals $\psi^{(l)}_C$, and let each unit of level $l+1$ decide through an aggregate $\psi^{(l+1)}_C=\sum_{C'\subset C}\mathsf w_{C'}\psi^{(l)}_{C'}$, such as a total demand. Then $r^{(l+1)}_{\mathrm{Hay}}=\dim\operatorname{span}\{\psi^{(l+1)}_C\}\le\min\{r^{(l)}_{\mathrm{Hay}},n^{(l+1)}\}$, with $n^{(l+1)}$ the number of units of level $l+1$: the public rank that suffices between meta-agents is nonincreasing up the hierarchy. A public message of that rank suffices for member $i$ of a unit only if $\psi_i\in\operatorname{span}\{\psi^{(l+1)}_C\}$, so a meta-agent $C$ whose top-down message must serve its members' own decisions needs rank $\dim\operatorname{span}\{\psi_i:i\in C\}$.
\end{proposition}
% src: H4-P14 (DERIVED); the reading after it is an interpretation
Proof in \cref{app:ren}. In our interpretation, hierarchy economizes communication between meta-agents exactly to the extent that decisions aggregate linearly, and economizes nothing top-down unless members' decisions are themselves functions of aggregates. Interpretation levels compound.

\begin{proposition}[Literal reading compounds across levels]\label{prop:ren-literal-levels}
Let every level run the literal update of \cref{prop:agy-literal-consensus}, the level-1 reading of posted precisions over noiseless channels.
\begin{enumerate}
\item[(a)] (Regular hierarchies.) Let $n_{\mathrm g}$ groups of $n$ members run literal consensus on connected $\Delta_1$-regular graphs, and one representative per group then on a connected $\Delta_2$-regular graph of groups. As the numbers of rounds at both levels grow, every representative's mean converges to $(N\Lambda_\circ m_\circ+\sum_i\Lambda^s_is_i)/(N\Lambda_\circ+\sum_i\Lambda^s_i)$, $N=n_{\mathrm g}n$: the common prior is counted $N$ times, as in flat literal consensus, and a regular hierarchy of group sizes $n_1,\dots,n_L$ counts it $n_1\cdots n_L=N$ times. Believed precisions diverge at every level.
\item[(b)] (Unequal groups.) After $n_{\mathrm c,1}$ rounds at level one, group $J$ enters level two with a weight proportional to $\rho(I+\mathrm{Adj}_J)^{n_{\mathrm c,1}}\mathsf v^{(J)}_{i_J}\mathsf v^{(J)\top}\Lambda_J(0)$ times its level-two centrality, with $\mathsf v^{(J)}$ the unit Perron vector of $I+\mathrm{Adj}_J$ and $i_J$ the representative; the group with the largest spectral radius dominates the top-level belief exponentially in $n_{\mathrm c,1}$, whatever its information, so the persuasion bias of \cref{prop:agy-literal-consensus} multiplies across levels.
\item[(c)] (Information form.) If every level runs consensus on information contributions (\cref{prop:agy-consensus}) and passes the exact sums $\sum_{i\in J}(h^s_i,\Lambda^s_i)$ upward, adding the common prior once at the top, the top-level belief is the union posterior for every hierarchy: the information form is closed under hierarchy, the posterior form is not.
\end{enumerate}
\end{proposition}
% src: H4-P15 (CHECKED); C3_09 recomputation; level-1 reading per HC_critic 2.8 (decision 10)
Proof in \cref{app:ren}; checked on four groups of five and on six groups of three (\cref{chk:ren-literal}). The overcounting is the overvaluation of \cref{sec:mis}; literal reading of top-down messages is treated in \cref{sec:hie}.

\subsection{Network renormalization}\label{sec:ren-network}

For networks without coordinates, the Laplacian renormalization group of \citet{villegas2023laplacian}, whose procedure we read in its primary text, builds the density matrix $\varrho_{\tau_{\mathrm L}}=e^{-\tau_{\mathrm L}L}/\tr e^{-\tau_{\mathrm L}L}$ of a graph Laplacian at a diffusion time $\tau_{\mathrm L}$ (not the correlation time $\tau$), locates characteristic scales at peaks of the specific heat of its normalized entropy, keeps the modes with eigenvalues strictly below $1/\tau_{\mathrm L}$ (a mode whose eigenvalue equals $1/\tau_{\mathrm L}$ is discarded, the convention at equality that we fix for every statement and computation below), forms as many macronodes as kept modes by joining node pairs in descending order of $\lvert(\varrho_{\tau_{\mathrm L}})_{ij}\rvert$, and rescales time; other network renormalizations \citep{garciaperez2018multiscale,garuccio2023multiscale,gabrielli2025network} are known to us from abstracts and a review. We apply the procedure to the coupling Laplacian $L^\Pi$ of the world-conditional precision.

\begin{proposition}[Laplacian renormalization does not see grounding]\label{prop:ren-lrg}
With uniform precision leak $\varepsilon^\Pi$,
\[
\frac{e^{-\tau_{\mathrm L}(\varepsilon^\Pi I+L^\Pi)}}{\tr e^{-\tau_{\mathrm L}(\varepsilon^\Pi I+L^\Pi)}}=\varrho_{\tau_{\mathrm L}},
\]
so the density matrix, the specific heat and the macronodes of the procedure do not depend on $\varepsilon^\Pi$; a cutoff applied to the spectrum of $\varepsilon^\Pi I+L^\Pi$ without subtracting its bottom removes every mode once $\tau_{\mathrm L}>1/\varepsilon^\Pi$, an artifact. Contrast depends on $\varepsilon^\Pi$: a region $C$ has $r(C)\ge r_0$ if and only if $\varepsilon^\Pi\le(\lambda_2(L_C)-r_0\cut(C)/\lvert C\rvert)/(r_0-1)$, which never holds if $\lambda_2(L_C)\le r_0\cut(C)/\lvert C\rvert$. A macronode is therefore selected at level $r_0$ exactly for leaks below this threshold, and it is a collective only if it also exceeds the baseline of its comparison medium.
\end{proposition}
% src: H4-P16 (CHECKED); baseline clause per decision 22

\begin{proposition}[Agreement for two equitable blocks]\label{prop:ren-lrg-agreement}
Let two blocks of equal size form a population whose cut is equitable, every member having cut coupling $\bar a^\partial$, with internal gaps at least $g_{\min}>2\bar a^\partial$. Then the spectrum of $L^\Pi$ is $0$, $2\bar a^\partial$ and the rest at least $g_{\min}$, and a renormalization step at any $\tau_{\mathrm L}$ with $2\bar a^\partial<1/\tau_{\mathrm L}\le\lambda_3(L^\Pi)$ and $\tau_{\mathrm L}(\lambda_3(L^\Pi)-2\bar a^\partial)>\ln N$ returns exactly the two blocks; such $\tau_{\mathrm L}$ exist when $g_{\min}/\bar a^\partial>2(1+\ln N)$, a condition on the attention graph alone. When both gaps equal $g$, the criterion $r\ge r_0$ selects each block if and only if $\varepsilon^\Pi\le(g-r_0\bar a^\partial)/(r_0-1)$ and the union if and only if $\varepsilon^\Pi\le2\bar a^\partial/(r_0-1)$, so the two procedures agree exactly for $2\bar a^\partial/(r_0-1)<\varepsilon^\Pi\le(g-r_0\bar a^\partial)/(r_0-1)$, an interval of leaks that contains the window of \cref{thm:col-bridge}(iv), and disagree outside it.
\end{proposition}
% src: H4-P17 (DERIVED); the union clause read as selection at a level, per H6-C3
Proofs in \cref{app:ren}. On the two grids of the bridge population of \cref{sec:col-bridge}, which are not equitable, the procedure returns the grids at $\tau_{\mathrm L}=11.2$ and $103$ and one macronode at $928$, as in the equitable case (\cref{chk:ren-lrg}). The union of two blocks that form a closed population is the whole population, never a collective under the exchangeable clause, on a torus or on an edge- and vertex-transitive declared graph (\cref{prop:col-whole-population}); there the criterion selects it only at levels below the baseline. With several scales a single step fails.

\begin{example}[Interleaved scales]\label{ex:ren-interleaved}
Let two stiff cliques $A$ and $A'$ of six agents (coupling $10$) be joined by a perfect matching of coupling $1$, each member of $A$ coupled to one member of $A'$, and let a soft $6\times6$ grid $B$ (coupling $1$) be attached to them by two edges of coupling $0.001$, one from each clique to one of two opposite corners of the grid. At $\varepsilon^\Pi=0.01$ the contrasts are $59.4$ for $A$ and $A'$, $197.7$ for $A\cup A'$ and $27.6$ for $B$, so the criterion $r\ge2$ selects all four sets. The Laplacian procedure returns $A$ and $A'$ as macronodes only for some $\tau_{\mathrm L}$ in $(0.0167,1)$, and not for $\tau_{\mathrm L}\in[0.283,0.306)$ or $[0.387,0.789)$; it returns $A\cup A'$ for $1\le\tau_{\mathrm L}<1/\lambda_2(L^\Pi)=4503$, and $B$ for $1/\lambda_3(L^\Pi)=3.73\le\tau_{\mathrm L}<4503$: no single diffusion time returns $A$, $A'$ and $B$ together, although every set that contrast selects appears somewhere along the flow.
\end{example}
% src: H4-X4 (CHECKED, H4_lrg.py); H4-E1 summarized after it
One step of the procedure uses one global timescale and fails whenever one collective's internal mixing time exceeds another's escape time, while contrast uses each region's own rates. On a homogeneous torus the procedure still forms macronodes, chained by single linkage into one cluster, where contrast selects nothing; in the nested medium it returns the blocks of each level, while contrast selects them according to the leak (\cref{chk:ren-lrg}; the return of the level blocks is proved for $\mathfrak q<b$ and every $\tau_{\mathrm L}$ in \cref{thm:obj-nested}(c)). Whether the sets of \cref{ex:ren-interleaved} are collectives depends on a comparison medium that we did not compute.

\begin{openproblem}[Which regions of high ratio Laplacian renormalization returns]\label{conj:ren-lrg}
With uniform leak, let $r_\star(C)=\lambda_2(L_C)\lvert C\rvert/\cut(C)$ be the grounding-free ratio of a region $C$. Characterize the regions $C$ that are macronodes of the Laplacian procedure at some $\tau_{\mathrm L}\in(1/\lambda_2(L_C),\lvert C\rvert/\cut(C))$. The conjecture that every region whose ratio $r_\star(C)$ exceeds a constant times $\ln N$ is one is false for every value of the constant (\cref{prop:obj-lrg-counter}).
\end{openproblem}
% src: H4-C1, withdrawn as a conjecture in the second repair pass: refuted as stated by prop:obj-lrg-counter (lane H); the restriction to local maximizers is contradicted numerically (ren_p2_lrg_restricted.py, chk:ren-lrg)
That conjecture is withdrawn as stated; its conclusion holds for two equitable blocks under the gap condition of \cref{prop:ren-lrg-agreement}, and a level-$l$ block of the nested medium with $\mathfrak q<b$ and $1\le l<n_{\mathrm H}$ is returned at some $\tau_{\mathrm L}$ in the interval exactly when $r_\star(C)>b/(b-\mathfrak q)$, a constant threshold without the factor $\ln N$ (\cref{thm:obj-nested}(c)). On the complete graph it fails for a region that omits one vertex, which is not a local maximizer of $r_\star$ under the addition or removal of one site. The restriction to such maximizers (\cref{open:obj-scale}(c)) is not supported by computation either: a clique of $16$ agents joined by one edge to a path of $256$ sites is a strict local maximizer with $r_\star=256$ and $\ln N=5.61$, and no scanned $\tau_{\mathrm L}$ in the interval returned it (\cref{chk:ren-lrg}), which is numerical evidence on a grid and not a proof. Laplacian renormalization asks at which diffusion scales the attention graph has structure; contrast asks which regions are collectives at the given grounding.

\subsection{Open problems}\label{sup:ren-open-problems}

\begin{openproblem}[Renormalization beyond the exact cases]\label{open:ren-beyond}
(a) Define contrast for the non-diffusive laws of block averages, which coarse-graining can produce (\cref{ex:ren-block-average}), extending \cref{open:col-non-diffusive}. (b) Prove discrete-time closure for non-reversible updates, beyond \cref{lem:met-discrete-closure}. (c) Determine the decay rate of decimated kernels in two and more dimensions (\cref{obs:ren-decay}). (d) Compute baselines of nested media embedded in lattices beyond one instance. (e) Resolve \cref{conj:ren-lrg} beyond two equitable blocks and the nested medium, and decide \cref{conj:sca-rank-down} for networks with many paths and heterogeneous metrics. (f) Propagate finite-round consensus errors across levels in the information form. (g) Bound the error of nonlinear and discrete renormalization, such as that of \cref{rem:ren-rgm}. (h) Prove the first-order rate of \cref{obs:ren-adiabatic-rate}.
\end{openproblem}
% src: H4 open items O2 to O8, O-H6-2, H1-O2, H1-O6
A partial result on the subject of this problem is in \cref{sec:apx}: for stable linear-Gaussian collectives in discrete time under small-gain coarse policies, \cref{cor:apx-uniform} bounds the relative entropy between the retained controlled path law and a declared Markov coarse model uniformly over the policy class, and \cref{thm:apx-levels} composes the defects through levels. Neither result settles any of the items (a) to (h).
